% Adapted from Spring 2018 Midterm 2 Q4b (Sequences)
% The original worked on the digits of integers; this version uses lists of digits.

\begin{blocksection}
\question [Adapted from CS 61A Spring 2018 Midterm 2 Q4b (Sequences)] Implement \texttt{combo}, which takes two lists of digits \texttt{a} and \texttt{b}. It returns the smallest list of digits that contains all of the digits of \texttt{a} in order, as well as all of the digits of \texttt{b} in order. A list is smaller than another if it is shorter, or if they have the same length and it comes first when compared digit by digit. You may use \texttt{smaller}, which returns the smaller of two lists.
\begin{lstlisting}
pub fn smaller(s: List(Int), t: List(Int)) -> List(Int) {
  // Implementation omitted
}

pub fn combo(a: List(Int), b: List(Int)) -> List(Int) {
  case a, b {
    [], _ -> ____
    _, [] -> a
    [x, ..ar], [y, ..br] if x == y -> ______________________
    [x, ..ar], [y, ..br] ->
      smaller(____________________, ____________________)
  }
}
assert combo([5, 3, 1], [4, 3, 2]) == [4, 5, 3, 1, 2]
assert combo([5, 3, 1], [4, 3, 2, 1]) == [4, 5, 3, 2, 1]
assert combo([1, 2, 3, 4], [9, 1, 2, 3]) == [9, 1, 2, 3, 4]
assert combo([], [3, 2, 1]) == [3, 2, 1]
\end{lstlisting}
\begin{solution}
\begin{lstlisting}
pub fn combo(a: List(Int), b: List(Int)) -> List(Int) {
  case a, b {
    [], _ -> b
    _, [] -> a
    [x, ..ar], [y, ..br] if x == y -> [x, ..combo(ar, br)]
    [x, ..ar], [y, ..br] ->
      smaller([x, ..combo(ar, b)], [y, ..combo(a, br)])
  }
}
\end{lstlisting}
\end{solution}
\begin{questionmeta}
%   The answer has to begin with either the first digit of \texttt{a} or the first digit of \texttt{b}; any other digit would be wasted. When the two first digits are equal, a single digit can serve both lists at once, so we use it once and continue with the rest of both lists.
%   When the first digits are different, we try both options. If the answer starts with \texttt{x}, that digit only accounts for the start of \texttt{a}, so we still need all of \texttt{b} afterward: \lstinline{combo(ar, b)}, not \lstinline{combo(ar, br)}. The other option is symmetric. \texttt{smaller} then picks the better of the two.
%   When one list is empty, the answer is just the other list, since every digit of it has to appear and nothing else is needed.
%   This is tree recursion with two branches, so it takes exponential time in the worst case.
%   The original question used integers like \texttt{531} instead of lists of digits, and it started from the last digit using \lstinline|% 10| and \lstinline|// 10|. Starting from the front of a list gives the same answers. In the integer version, ``smallest'' compared numbers, which is why \texttt{smaller} first compares lengths and then digits.
\end{questionmeta}
\end{blocksection}
